% !TEX program = pdflatex
\documentclass[11pt]{article}

% ==========================
% Encoding / language (pdflatex)
% ==========================
\usepackage[utf8]{inputenc}
\usepackage[T1]{fontenc}
\usepackage[english]{babel}
\usepackage{lmodern}
\usepackage{microtype}

% ==========================
% Layout
% ==========================
\usepackage[margin=1in]{geometry}

% ==========================
% Math
% ==========================
\usepackage{amsmath,amssymb,amsthm,mathtools}

% ==========================
% Lists / links
% ==========================
\usepackage{enumitem}
\usepackage[hidelinks]{hyperref}

% ==========================
% Macros
% ==========================
\newcommand{\Pbb}{\mathbb{P}}
\newcommand{\E}{\mathbb{E}}
\newcommand{\R}{\mathbb{R}}
\newcommand{\F}{\mathcal{F}}
\newcommand{\1}{\mathbf{1}}
\newcommand{\dd}{\,\mathrm{d}}

% ==========================
% Theorem environments
% ==========================
\numberwithin{equation}{section}

\theoremstyle{plain}
\newtheorem{theorem}{Theorem}[section]
\newtheorem{lemma}[theorem]{Lemma}
\newtheorem{proposition}[theorem]{Proposition}
\newtheorem{corollary}[theorem]{Corollary}

\theoremstyle{definition}
\newtheorem{definition}[theorem]{Definition}

\theoremstyle{remark}
\newtheorem{remark}[theorem]{Remark}

% ==========================
% Board environment (robust)
% ==========================
\newenvironment{board}[1][On the board (plan + comments)]{%
  \par\medskip
  \noindent
  \fbox{%
    \begin{minipage}{0.98\linewidth}
    \textbf{#1.}\par\medskip
}{%
    \end{minipage}%
  }%
  \par\medskip
}

% ==========================
% Title
% ==========================
\title{Lecture Notes (Board Plan + Proofs) --- It\^o Calculus and It\^o's Formula}
\author{}
\date{}

\begin{document}
\maketitle
\vspace{-1em}

\noindent\textbf{Goal of this lecture.}
We explain why classical Riemann--Stieltjes integration fails for Brownian motion,
introduce quadratic variation, and derive It\^o's formula (1D and multi-D).
We then work through standard patterns (powers, exponentials, logarithms),
solve GBM and OU/Vasicek, and connect to finance via discounting and self-financing gains.

\textbf{To do on board}
\begin{itemize}[leftmargin=2em]
  \item Motivation: Brownian paths are too rough $\Rightarrow$ usual calculus breaks.
  \item Key new object: quadratic variation $[W]_t = t$.
  \item It\^o formula: Taylor expansion + the rule $(\Delta W)^2 \approx \Delta t$.
  \item Finance: discounting to force martingales, and self-financing gains $\int \phi\,\dd S$.
\end{itemize}


% ==========================================================
\section{Warm-up recap: It\^o integral as an $L^2$ limit}
% ==========================================================

\subsection{Construction on step processes}

\begin{definition}[It\^o integral for step processes]
Let $\phi$ be a step (simple adapted) process on $[0,T]$:
\[
\phi_t=\sum_{k=0}^{n-1}\phi_{t_k}\,\1_{(t_k,t_{k+1}]}(t),
\qquad \phi_{t_k}\in\F_{t_k}.
\]
Define
\[
\int_0^T \phi_t\,\dd W_t
:=\sum_{k=0}^{n-1}\phi_{t_k}\big(W_{t_{k+1}}-W_{t_k}\big).
\]
\end{definition}

\textbf{To do on board}
\begin{itemize}[leftmargin=2em]
  \item Stress: this is a \emph{definition} on step processes.
  \item Left point = non-anticipative: $\phi_{t_k}$ chosen before $\Delta W_k$ is seen.
\end{itemize}


\subsection{Two core computations}

\begin{proposition}[Mean zero for step integrands]
If $\phi$ is a step adapted process, then
\[
\E\!\left[\int_0^T \phi_t\,\dd W_t\right]=0.
\]
\end{proposition}

\begin{proof}
Write $I=\sum_{k}\phi_{t_k}\Delta W_k$ with $\Delta W_k:=W_{t_{k+1}}-W_{t_k}$.
Then
\[
\E[\phi_{t_k}\Delta W_k]=\E\!\big[\E[\phi_{t_k}\Delta W_k\mid \F_{t_k}]\big].
\]
Since $\phi_{t_k}$ is $\F_{t_k}$-measurable and $\Delta W_k$ is independent of $\F_{t_k}$ with mean $0$,
\[
\E[\phi_{t_k}\Delta W_k\mid \F_{t_k}]
=\phi_{t_k}\E[\Delta W_k]=0.
\]
Sum over $k$.
\end{proof}

\begin{theorem}[It\^o isometry for step integrands]
If $\phi$ is a step adapted process, then
\[
\E\!\left[\left(\int_0^T \phi_t\,\dd W_t\right)^2\right]
=\E\!\left[\int_0^T \phi_t^2\,\dd t\right].
\]
\end{theorem}

\begin{proof}
Let $I=\sum_k \phi_{t_k}\Delta W_k$.
Expand:
\[
I^2=\sum_k \phi_{t_k}^2(\Delta W_k)^2 + 2\sum_{i<j}\phi_{t_i}\phi_{t_j}\Delta W_i\Delta W_j.
\]
For the diagonal terms, condition on $\F_{t_k}$:
\[
\E[\phi_{t_k}^2(\Delta W_k)^2]
=\E\!\big[\phi_{t_k}^2\E[(\Delta W_k)^2]\big]
=\E[\phi_{t_k}^2](t_{k+1}-t_k).
\]
For cross terms $i<j$, condition on $\F_{t_j}$: then
$\phi_{t_i},\phi_{t_j},\Delta W_i$ are $\F_{t_j}$-measurable while
$\Delta W_j$ is independent of $\F_{t_j}$ with mean $0$, so the conditional expectation is $0$.
Hence cross terms vanish, giving
\[
\E[I^2]=\sum_k \E[\phi_{t_k}^2](t_{k+1}-t_k)=\E\!\left[\int_0^T \phi_t^2\,\dd t\right].
\]
\end{proof}

\textbf{To do on board}[On the board (the “one big computation”)]
\begin{itemize}[leftmargin=2em]
  \item Expand the square, then show \emph{cross terms vanish}.
  \item Key mechanism: independent increments + mean zero.
\end{itemize}


\subsection{Extension to general integrands}

\begin{definition}[It\^o integral in $L^2$]
Let $\phi$ be adapted and satisfy $\E\int_0^T \phi_t^2\,\dd t<\infty$.
Choose step processes $\phi^n$ such that $\phi^n\to\phi$ in $L^2(\Omega\times[0,T])$.
Define
\[
\int_0^T \phi_t\,\dd W_t := \lim_{n\to\infty}\int_0^T \phi_t^n\,\dd W_t
\quad \text{in } L^2(\Omega).
\]
\end{definition}

\begin{proposition}[Well-definedness]
The limit does not depend on the choice of the approximating sequence $(\phi^n)$.
\end{proposition}

\begin{proof}
If $\phi^n\to\phi$ and $\psi^n\to\phi$ in $L^2(\Omega\times[0,T])$, then $\phi^n-\psi^n\to 0$.
By the isometry (applied to step processes),
\[
\E\!\left[\left(\int_0^T(\phi_t^n-\psi_t^n)\,\dd W_t\right)^2\right]
=\E\!\left[\int_0^T(\phi_t^n-\psi_t^n)^2\,\dd t\right]\to 0.
\]
So the two limits in $L^2(\Omega)$ coincide.
\end{proof}

% ==========================================================
\section{Quadratic variation and why Riemann--Stieltjes fails}
% ==========================================================

\subsection{Quadratic variation of Brownian motion}

\begin{definition}[Quadratic variation along a partition]
For a partition $\pi=\{0=t_0<\cdots<t_n=t\}$ define
\[
[W]_t^{\pi}:=\sum_{k=0}^{n-1}\big(W_{t_{k+1}}-W_{t_k}\big)^2.
\]
\end{definition}

\begin{theorem}[Quadratic variation: uniform partitions]
For the uniform grid $t_k=kt/n$,
\[
\sum_{k=0}^{n-1}(W_{t_{k+1}}-W_{t_k})^2 \xrightarrow[n\to\infty]{L^2} t.
\]
\end{theorem}

\begin{proof}
Let $\Delta W_k:=W_{t_{k+1}}-W_{t_k}$; then $\Delta W_k\sim\mathcal{N}(0,\Delta t)$ i.i.d.,
with $\Delta t=t/n$.

\smallskip\noindent
\textbf{Step 1: expectation.}
\[
\E\!\left[\sum_k(\Delta W_k)^2\right]
=\sum_k \E[(\Delta W_k)^2]
=\sum_k \Delta t = n\Delta t=t.
\]

\smallskip\noindent
\textbf{Step 2: variance.}
Independence gives
\[
\mathrm{Var}\!\left(\sum_k(\Delta W_k)^2\right)=\sum_k \mathrm{Var}((\Delta W_k)^2).
\]
If $Z\sim\mathcal{N}(0,\sigma^2)$, then $\E[Z^4]=3\sigma^4$, so
$\mathrm{Var}(Z^2)=3\sigma^4-\sigma^4=2\sigma^4$.
Here $\sigma^2=\Delta t$, so $\mathrm{Var}((\Delta W_k)^2)=2(\Delta t)^2$.
Hence
\[
\mathrm{Var}\!\left(\sum_k(\Delta W_k)^2\right)=n\cdot 2(\Delta t)^2
=2n\left(\frac{t}{n}\right)^2=\frac{2t^2}{n}\to 0.
\]
Thus $L^2$ convergence holds.
\end{proof}

\textbf{To do on board}
\begin{itemize}[leftmargin=2em]
  \item Write the expectation and variance computation cleanly.
  \item Message: for Brownian motion, the \emph{squares} accumulate $\Rightarrow [W]_t=t$.
\end{itemize}


\subsection{Why the sampling point matters}

Let $0=t_0<\cdots<t_n=T$ be uniform and $\Delta W_k=W_{t_{k+1}}-W_{t_k}$.
Define
\[
L_n:=\sum_{k=0}^{n-1}W_{t_k}\Delta W_k,
\qquad
R_n:=\sum_{k=0}^{n-1}W_{t_{k+1}}\Delta W_k.
\]

\begin{proposition}[Right minus left equals quadratic variation]
\[
R_n-L_n=\sum_{k=0}^{n-1}(\Delta W_k)^2.
\]
\end{proposition}

\begin{proof}
Simply expand:
\[
R_n-L_n=\sum_k (W_{t_{k+1}}-W_{t_k})\Delta W_k=\sum_k (\Delta W_k)^2.
\]
\end{proof}

By quadratic variation, $\sum_k(\Delta W_k)^2\to T$ (in probability and even in $L^2$ for uniform grids).
Hence \emph{if} $L_n$ converges, then $R_n$ must converge to a limit shifted by $T$.

\textbf{To do on board}
\begin{itemize}[leftmargin=2em]
  \item “Changing the evaluation point changes the limit” when integrand is as rough as $W$.
  \item This is the conceptual reason Riemann--Stieltjes is not robust here.
\end{itemize}


\subsection{Canonical identity: $\int_0^T W_t\,\dd W_t$}

\begin{proposition}[Key identity]
\[
\int_0^T W_t\,\dd W_t=\frac12\big(W_T^2-T\big).
\]
\end{proposition}

\begin{proof}
Use the algebraic identity
\[
W_{t_{k+1}}^2-W_{t_k}^2 = 2W_{t_k}\Delta W_k + (\Delta W_k)^2.
\]
Sum over $k$:
\[
W_T^2-W_0^2=2\sum_k W_{t_k}\Delta W_k + \sum_k(\Delta W_k)^2.
\]
Thus
\[
\sum_k W_{t_k}\Delta W_k
=\frac12\Big(W_T^2-\sum_k(\Delta W_k)^2\Big)\xrightarrow[n\to\infty]{}\frac12(W_T^2-T),
\]
since $\sum_k(\Delta W_k)^2\to T$.
By definition, the limit of left sums is $\int_0^T W_t\,\dd W_t$.
\end{proof}

\textbf{To do on board}
\begin{itemize}[leftmargin=2em]
  \item Extremely accessible proof: algebra + $[W]_T=T$.
  \item Highlight the appearance of the correction term $-\tfrac12T$.
\end{itemize}


% ==========================================================
\section{It\^o processes and It\^o's formula (1D)}
% ==========================================================

\subsection{It\^o processes}

\begin{definition}[It\^o process]
A real-valued adapted process $X$ is an It\^o process if
\[
X_t=X_0+\int_0^t a_s\,\dd s+\int_0^t b_s\,\dd W_s,
\]
where $a$ is adapted with $\E\int_0^T|a_s|\,\dd s<\infty$ and
$b$ is adapted with $\E\int_0^T b_s^2\,\dd s<\infty$.
We write in differential form: $\dd X_t=a_t\,\dd t+b_t\,\dd W_t$.
\end{definition}

\textbf{To do on board}
\begin{itemize}[leftmargin=2em]
  \item Emphasize: drift ($\dd t$) + diffusion ($\dd W$).
  \item Integrability conditions guarantee the integrals are well-defined.
\end{itemize}


\subsection{It\^o's formula (scalar case)}

\begin{theorem}[It\^o's formula in 1D]
Let $\dd X_t=a_t\,\dd t+b_t\,\dd W_t$ and $f\in C^{1,2}([0,T]\times\R)$.
Then $Y_t=f(t,X_t)$ satisfies
\[
\dd f(t,X_t)
=\Big(\partial_t f + a_t\,\partial_x f + \tfrac12 b_t^2\,\partial_{xx} f\Big)(t,X_t)\,\dd t
+\big(b_t\,\partial_x f\big)(t,X_t)\,\dd W_t.
\]
\end{theorem}

\begin{proof}[Proof idea (what you can present on the board)]
Fix a partition $0=t_0<\cdots<t_n=t$ with mesh $|\pi|\to 0$.
Write the increment
\[
f(t_{k+1},X_{t_{k+1}})-f(t_k,X_{t_k})
=
\Big[f(t_{k+1},X_{t_{k+1}})-f(t_k,X_{t_{k+1}})\Big]
+
\Big[f(t_k,X_{t_{k+1}})-f(t_k,X_{t_k})\Big].
\]
Apply Taylor:
\[
f(t_{k+1},x)-f(t_k,x)=\partial_t f(t_k,x)\Delta t_k + o(\Delta t_k),
\]
and
\[
f(t_k,X_{t_{k+1}})-f(t_k,X_{t_k})
=\partial_x f(t_k,X_{t_k})\Delta X_k + \tfrac12 \partial_{xx} f(t_k,X_{t_k})(\Delta X_k)^2 + r_k,
\]
where the remainder $r_k$ is of higher order (controlled using smoothness of $f$ and moment bounds).

Now use $\Delta X_k = a_{t_k}\Delta t_k + b_{t_k}\Delta W_k + \text{smaller terms}$.
The key scaling is:
\[
\Delta W_k = O(\sqrt{\Delta t_k}),
\qquad
(\Delta W_k)^2 \approx \Delta t_k,
\qquad
\Delta t_k\Delta W_k \approx 0,
\qquad
(\Delta t_k)^2 \approx 0.
\]
Thus
\[
(\Delta X_k)^2 = b_{t_k}^2(\Delta W_k)^2 + \text{terms negligible after summation}.
\]
Summing over $k$ and letting $|\pi|\to 0$, the sums converge (in probability, and in $L^2$ under standard assumptions) to
\[
\int_0^t \partial_t f(s,X_s)\,\dd s
+\int_0^t \partial_x f(s,X_s)\,\dd X_s
+\frac12\int_0^t \partial_{xx} f(s,X_s)\,b_s^2\,\dd s,
\]
which yields the formula.
\end{proof}

\textbf{To do on board}[On the board (It\^o formula derivation)]
\begin{itemize}[leftmargin=2em]
  \item Start from Taylor expansion up to second order in space.
  \item Insert $\Delta X \approx a\Delta t + b\Delta W$.
  \item Use the “multiplication table”: $(\Delta W)^2\approx \Delta t$, $\Delta t\,\Delta W\approx 0$.
  \item Conclude: extra drift term $\tfrac12 b^2 f_{xx}\,\dd t$.
\end{itemize}


\subsection{Quadratic variation reminder for It\^o processes}
If $\dd X_t=a_t\,\dd t+b_t\,\dd W_t$, then
\[
[X]_t=\int_0^t b_s^2\,\dd s,
\qquad\text{so}\qquad \dd[X]_t=b_t^2\,\dd t.
\]
This is consistent with the heuristic $(\dd W_t)^2=\dd t$.

% ==========================================================
\section{Multidimensional It\^o formula + product rule}
% ==========================================================

\subsection{Quadratic covariation}
Let $X\in\R^m$ satisfy $\dd X_t=a_t\,\dd t + B_t\,\dd W_t$ where $W\in\R^d$.
If $W$ has covariance matrix $\Sigma$ (i.e.\ $\langle W^i,W^j\rangle_t=\rho_{ij}t$),
then
\[
\dd[X]_t = B_t\,\Sigma\,B_t^\top\,\dd t.
\]

\begin{theorem}[Multidimensional It\^o formula]
Let $f\in C^{1,2}([0,T]\times\R^m)$ and $Y_t=f(t,X_t)$. Then
\[
\dd f(t,X_t)
=
\partial_t f(t,X_t)\,\dd t
+
\nabla f(t,X_t)^\top \dd X_t
+
\frac12 \mathrm{Tr}\!\big(\nabla^2 f(t,X_t)\,\dd[X]_t\big).
\]
Equivalently,
\[
\dd f(t,X_t)
=
\Big(\partial_t f+\nabla f^\top a_t+\tfrac12\mathrm{Tr}(\nabla^2 f\,B_t\Sigma B_t^\top)\Big)(t,X_t)\,\dd t
+
\nabla f(t,X_t)^\top B_t\,\dd W_t.
\]
\end{theorem}

\textbf{To do on board}
\begin{itemize}[leftmargin=2em]
  \item State clearly: gradient term + Hessian trace term.
  \item Explain correlations appear through $\Sigma$ (cross terms in quadratic covariation).
\end{itemize}


\subsection{Product rule}
Take $f(x,y)=xy$ and apply the 2D It\^o formula:
\[
\dd(X_tY_t)=X_t\,\dd Y_t + Y_t\,\dd X_t + \dd[X,Y]_t.
\]
For Brownian motion, $\dd[W,W]_t=\dd t$, giving $\dd(W_t^2)=2W_t\,\dd W_t+\dd t$.

% ==========================================================
\section{Standard patterns (powers, exponentials, logs)}
% ==========================================================

\subsection{Powers: $W_t^n$}
For $f(x)=x^n$, $f'(x)=n x^{n-1}$, $f''(x)=n(n-1)x^{n-2}$, hence
\[
\dd(W_t^n)=nW_t^{n-1}\,\dd W_t + \frac12 n(n-1)W_t^{n-2}\,\dd t.
\]
Taking expectations:
\[
\frac{\dd}{\dd t}\E[W_t^n]=\frac12 n(n-1)\E[W_t^{n-2}].
\]

\subsection{Exponentials: exponential martingales}
For $Y_t=e^{\lambda W_t}$,
\[
\dd(e^{\lambda W_t})=\lambda e^{\lambda W_t}\,\dd W_t + \frac12\lambda^2 e^{\lambda W_t}\,\dd t.
\]
Thus the corrected process
\[
M_t:=\exp\!\Big(\lambda W_t-\tfrac12\lambda^2 t\Big)
\]
satisfies $\dd M_t=\lambda M_t\,\dd W_t$, so $M$ is a (local) martingale (in fact a true martingale for fixed $t$).

\subsection{Logarithms: $\ln S_t$ for $S_t>0$}
If $\dd S_t=\mu_t\,\dd t+\sigma_t\,\dd W_t$ and $S_t>0$, then for $f(x)=\ln x$:
\[
\dd(\ln S_t)=\frac{1}{S_t}\,\dd S_t-\frac12\frac{1}{S_t^2}\,\dd[S]_t
=
\Big(\frac{\mu_t}{S_t}-\frac12\frac{\sigma_t^2}{S_t^2}\Big)\dd t+\frac{\sigma_t}{S_t}\dd W_t.
\]
If $\dd S_t=\mu_t S_t\,\dd t+\sigma_t S_t\,\dd W_t$ (GBM form), this simplifies to
$\dd\ln S_t=(\mu_t-\tfrac12\sigma_t^2)\dd t+\sigma_t\dd W_t$.

% ==========================================================
\section{Solving canonical SDEs: GBM and OU/Vasicek}
% ==========================================================

\subsection{GBM solution}
If $\dd S_t=\mu S_t\,\dd t+\sigma S_t\,\dd W_t$ with $S_0>0$, then
\[
\dd(\ln S_t)=\Big(\mu-\tfrac12\sigma^2\Big)\dd t+\sigma\,\dd W_t,
\]
so integrating:
\[
\ln S_t=\ln S_0+\Big(\mu-\tfrac12\sigma^2\Big)t+\sigma W_t,
\qquad
S_t=S_0\exp\!\Big(\big(\mu-\tfrac12\sigma^2\big)t+\sigma W_t\Big).
\]
Hence $\ln S_t$ is Gaussian and $S_t$ is lognormal.

\textbf{To do on board}
\begin{itemize}[leftmargin=2em]
  \item Compute $\dd\ln S_t$ once carefully (students must master this).
  \item Then integrate and conclude distribution immediately.
\end{itemize}


\subsection{OU/Vasicek}
For $\dd r_t=k(m-r_t)\dd t+\sigma \dd W_t$ with $k>0$, set $Y_t=e^{kt}r_t$.
Then product rule gives
\[
\dd(e^{kt}r_t)=ke^{kt}r_t\,\dd t+e^{kt}\dd r_t
=km e^{kt}\,\dd t+\sigma e^{kt}\,\dd W_t.
\]
Integrate:
\[
r_t=m+(r_0-m)e^{-kt}+\sigma\int_0^t e^{-k(t-s)}\,\dd W_s.
\]

% ==========================================================
\section{Finance: discounting and self-financing gains}
% ==========================================================

\subsection{Discounting trick}
Assume $\dd S_t=\mu_t S_t\,\dd t+\sigma_t S_t\,\dd W_t$.
Let
\[
A_t:=\exp\!\Big(-\int_0^t \gamma_s\,\dd s\Big),\qquad X_t:=A_t S_t.
\]
Since $A$ has finite variation, $\dd[A,S]_t=0$ and $\dd A_t=-\gamma_t A_t\,\dd t$.
Thus
\[
\dd X_t = A_t\,\dd S_t + S_t\,\dd A_t
= A_t S_t(\mu_t-\gamma_t)\dd t + A_t S_t\sigma_t\,\dd W_t.
\]
Choose $\gamma_t=\mu_t$ to kill the drift:
\[
\dd X_t = X_t\sigma_t\,\dd W_t,
\]
so $X$ is a (local) martingale.

\textbf{To do on board}
\begin{itemize}[leftmargin=2em]
  \item Compute $\dd(A_tS_t)$ and emphasize: finite variation $\Rightarrow \dd[A,S]=0$.
  \item Interpretation: discount by the “right drift” to get a martingale.
\end{itemize}


\subsection{Self-financing gains: discrete $\to$ continuous}
On a trading grid $0=t_0<\cdots<t_n=T$, holding $\phi_{t_k}$ shares on $(t_k,t_{k+1}]$ yields gains
\[
G_T^\pi(\phi)=\sum_{k=0}^{n-1}\phi_{t_k}(S_{t_{k+1}}-S_{t_k}).
\]
As mesh $\to 0$, this converges to $\int_0^T \phi_t\,\dd S_t$ (It\^o integral),
and we write heuristically $\dd V_t=\phi_t\,\dd S_t$ for a self-financing portfolio.

\textbf{To do on board}
\begin{itemize}[leftmargin=2em]
  \item Stress predictability: $\phi_{t_k}$ is chosen using information at $t_k$.
  \item This is the finance reason behind the left-point convention.
\end{itemize}


% ==========================================================
\section{Common mistakes (quick checklist)}
% ==========================================================

\begin{remark}[The 5 classic traps]
\begin{enumerate}[leftmargin=2em]
  \item Forgetting the It\^o correction term $\tfrac12 f_{xx}\,b^2\,\dd t$.
  \item Writing $(\dd W_t)^2=0$ (wrong: $(\dd W_t)^2=\dd t$).
  \item Mixing variables: $x$ vs $X_t$ in $\partial_x f(t,x)$.
  \item Using the usual chain rule as if paths were differentiable.
  \item Using right-point sums (anticipating) when the model is non-anticipative.
\end{enumerate}
\end{remark}

\textbf{To do on board}[Final message to leave on the board]
\[
(\dd W_t)^2=\dd t,\qquad \dd t\,\dd W_t=0,\qquad (\dd t)^2=0,
\]
\[
\dd f(t,X_t)=f_t\,\dd t+f_x\,\dd X_t+\tfrac12 f_{xx}\,\dd[X]_t.
\]


\end{document}